\documentclass{report}
\usepackage{amsmath,amssymb}
\setlength{\parindent}{0mm}
\setlength{\parskip}{1em}
\begin{document}
\begin{center}
\rule{6in}{1pt} \
{\large Martin Schulz \\
{\bf Addressing Performance Challenges on Modern Node Architectures}}

PO Box 808 \\ L-561 \\ Livermore \\ CA 94551 \\ USA
\\
{\tt schulzm@llnl.gov}\end{center}

Modern HPC systems feature complex node architectures with a large number
of cores, deep memory hierarchies, non-uniform memory access (NUMA)
properties and hyper threading; and future systems will only grow in
complexity. When porting solver libraries to such systems, understanding
the impact of data access patterns is key to extracting good performance
as well as energy efficiency. In this talk I will present a set of tools,
which can help developers extract memory access patterns at runtime,
annotate it with the necessary context to allow root cause analysis and
visualize the results in an intuitive way.

Mitos is a sampling tool that can transparently capture memory accesses
at runtime with minimal overhead at runtime; Caliper is an annotation
library that can be used to add information about a code�s structure,
such as key data structures or phases, in a composable and reusable
manner; and MemAxes is a visualization tool that maps memory access
information, combined with structural information, back to hardware
topologies, key data structures, or even the underlying matrices or
simulation domains. Combined these tools provide the necessary insight
into an application�s or a library�s memory access behavior to optimize
data communication for complex node architectures.


\end{document}
